\begin{itemframe}{الگوریتم‌های برنامه ریزی خطی}
\itm
الگوریتم‌های متعددی برای حل برنامه‌های خطی توسعه یافته‌اند که برخی از آن‌ها در زمان چندجمله‌ای اجرا می‌شوند و برخی دیگر نه، اما همگی آن‌ها بسیار پیچیده‌اند.
\itm
بنابراین، در اینجا این الگوریتم‌ها را نشان نمی‌دهیم. در عوض مثالی خواهیم دید که برخی از ایده‌های پشت الگوریتم سیمپلکس
\fn{simplex algorithm}
را نشان می‌دهد؛ الگوریتمی که در حال حاضر رایج‌ترین روش حل است.

\end{itemframe}


\begin{itemframe-s}{الگوریتم‌های برنامه ریزی خطی}{برنامه‌های خطی در حالت کلی}
\itm
در مسئله برنامه‌ریزی خطی، هدف بهینه‌سازی یک تابع خطی با رعایت مجموعه‌ای از نامعادلات خطی است.
\itm
فرض کنید اعداد حقیقی
$a_1, a_2, \dots, a_n$
 و متغیرهای
$x_1, x_2, \dots, x_n$
داده شده باشند. آنگاه تابع خطی
\fn{linear function}
 $f$
روی این متغیرها به صورت زیر تعریف می‌شود:
$$
f(x_1, x_2, \dots, x_n) = a_1 x_1 + a_2 x_2 + \cdots + a_n x_n = \sum_{j=1}^{n} a_j x_j
$$
\end{itemframe-s}


\begin{itemframe-s}{الگوریتم‌های برنامه ریزی خطی}{برنامه‌های خطی در حالت کلی}
\itm
حال اگر $b$ عددی حقیقی و $f$ تابعی خطی باشد، آنگاه معادله
$$
f(x_1, x_2, \dots, x_n) = b
$$
یک تساوی خطی است و نامعادلات
$$
f(x_1, x_2, \dots, x_n) \ge b \quad \text{و} \quad f(x_1, x_2, \dots, x_n) \le b
$$
نیز نامعادلات خطی هستند. ما اصطلاح کلی «محدودیت خطی»
\fn{linear constraints}
را برای اشاره به هر دو نوع به کار می‌بریم.
\itm
توجه داشته باشید در برنامه‌ریزی خطی نامساوی اکید
\fn{strict inequalitie}
را مجاز نیست. (
$<, >$
)
\end{itemframe-s}


\begin{itemframe-s}{الگوریتم‌های برنامه ریزی خطی}{برنامه‌های خطی در حالت کلی}
\itm
به‌طور رسمی، مسئله برنامه‌ریزی خطی به معنای کمینه‌سازی یا بیشینه‌سازی یک تابع خطی تحت مجموعه‌ای متناهی از محدودیت‌های خطی است. اگر هدف کمینه‌سازی باشد، آن را «برنامه خطی کمینه‌سازی»
\fn{minimization linear program}
و در غیر این صورت «برنامه خطی بیشینه‌سازی»
\fn{maximization linear program}
می‌نامیم.
\itm
برای بحث درباره الگوریتم‌ها و ویژگی‌های برنامه‌ریزی خطی، بهتر است که از نمادگذاری استاندارد برای ورودی‌ها استفاده کنیم.
\itm
به‌طور قراردادی، یک برنامه خطی بیشینه‌سازی شامل ورودی‌هایی به‌صورت $n$ عدد حقیقی
$c_1, c_2, \dots, c_n$،
$m$
عدد حقیقی
$b_1, b_2, \dots, b_m$
و
$mn$
عدد حقیقی
$a_{ij}$
به ازای
$i = 1, 2, \dots, m$
و
$j = 1, 2, \dots, n$
است.
\end{itemframe-s}


\begin{itemframe-s}{الگوریتم‌های برنامه ریزی خطی}{برنامه‌های خطی در حالت کلی}
\itm
هدف یافتن $n$ عدد حقیقی
$x_1, x_2, \dots, x_n$
است که عبارت زیر (تابع هدف
\fn{objective function}
) را بیشینه کند:
$$
\sum_{j=1}^{n} c_j x_j
$$

با رعایت محدودیت‌های:

$$
\sum_{j=1}^{n} a_{ij} x_j \ge b_i \quad for \:  i = 1, 2, \dots, m
$$

$$
x_j \ge 0 \quad for \: j = 1, 2, \dots, n
$$
\end{itemframe-s}


\begin{itemframe-s}{الگوریتم‌های برنامه ریزی خطی}{برنامه‌های خطی در حالت کلی}
\itm
برای فشرده‌تر کردن این بیان، می‌توانیم ماتریس
$A = (a_{ij})$
از ابعاد
$m \times n$،
بردار $b = (b_i)$ از اندازه $m$،
بردار $c = (c_j)$ از اندازه $n$،
و بردار $x = (x_j)$ از اندازه $n$ را تعریف کنیم.
آنگاه برنامه خطی به‌صورت زیر بازنویسی می‌شود:
$$
\text{  maximize} c^T x
$$

$$
\text{  to subject} Ax \ge b
$$

$$
x \ge 0
$$
\end{itemframe-s}


\begin{itemframe-s}{الگوریتم‌های برنامه ریزی خطی}{برنامه‌های خطی در حالت کلی}
\itm
درک معادلات بالا نیاز به کمی دانش جبر خطی دارد؛
%todo do we need LA in appendix.
$c^T x$
ضرب داخلی دو بردار $n$ تایی است.
$Ax$
حاصل‌ضرب ماتریس
$m \times n$
در بردار $n$تایی، و
$x \ge 0$
بدین معنی است که تمامی مؤلفه‌های بردار $x$ باید نامنفی باشند.
\itm
این نوع نمایش را «فرم استاندارد»
\fn{standard form}
یک برنامه خطی می‌نامیم، و در سراسر این فصل فرض می‌کنیم که $A$، $b$، و $c$ دارای ابعاد فوق هستند.
\itm
فرم استاندارد لزوماً با موقعیت‌های واقعی که می‌خواهید مدل‌سازی کنید مطابقت ندارد. برای مثال، ممکن است محدودیت‌های تساوی داشته باشید یا متغیرهایی که می‌توانند مقادیر منفی بگیرند. اما می‌توان هر برنامه خطی را به این فرم استاندارد تبدیل کرد.

\end{itemframe-s}


\begin{itemframe-s}{الگوریتم‌های برنامه ریزی خطی}{برنامه‌های خطی در حالت کلی}
\itm
در ادامه اصطلاحاتی را برای توصیف راه‌حل‌های برنامه‌های خطی معرفی می‌کنیم. یک مقداردهی خاص به متغیر
$x$
را با نماد
$\bar{x}$
نشان می‌دهیم. اگر
$\bar{x}$
تمام محدودیت‌ها را ارضا کند، آنگاه یک راه‌حل مجاز
\fn{feasible solution}
است. در غیر این صورت، یک راه‌حل غیرمجاز
\fn{infeasible solution}
است.
\itm
یک راه‌حل مانند
$\bar{x}$
دارای مقدار هدف
\fn{objective value}
$c^T \bar{x}$
است. اگر
$\bar{x}$
یک راه‌حل «مجاز» باشد و مقدار هدف آن در میان همه راه‌حل‌های مجاز بیشینه باشد، آنگاه
$\bar{x}$
یک «راه‌حل بهینه»
\fn{optimal solution}
است و
$c^T \bar{x}$
مقدار هدف بهینه
\fn{optimal objective value}
نامیده می‌شود.
\end{itemframe-s}


\begin{itemframe-s}{الگوریتم‌های برنامه ریزی خطی}{برنامه‌های خطی در حالت کلی}
\itm
اگر یک برنامه خطی هیچ راه‌حل مجازی نداشته باشد، آن را یک برنامه خطی ناممکن
\fn{infeasible}
و در غیر این صورت آن را یک برنامه خطی ممکن
\fn{feasible}
می‌نامیم.
\itm
اگر یک دستگاه مختصات با n محور فرض کنیم، آنگاه هر نقطه روی این دستگاه یک مقدار دهی به بردار $x$ است. ناحیه‌ایی روی این نمودار را که همه نقاط آن پاسخ‌های مجاز هستند «ناحیه مجاز»
\fn{feasible region}
می‌نامیم.

\itm
اگر یک برنامه خطی راه‌حل مجاز داشته باشد، اما «مقدار هدف بهینهٔ» متناهی نداشته باشد، آنگاه ناحیه مجاز و در نتیجه برنامه خطی «نامحدود»
\fn{unbounded}
است. عکس این قضیه الزاماً صادق نیست.
\end{itemframe-s}


\begin{itemframe-s}{الگوریتم‌های برنامه ریزی خطی}{برنامه‌های خطی در حالت کلی}
\itm
به‌طور کلی، برنامه‌های خطی را می‌توان به شکل کارآمد حل کرد.
دو دسته الگوریتم، به نام‌های الگوریتم‌های بیضوی
\fn{ellipsoid algorithms}
و الگوریتم‌های نقطه-داخلی
\fn{interior-point algorithms}
، وجود دارند که برنامه‌های خطی را در زمان چندجمله‌ای حل می‌کنند.
\itm
افزون بر این‌ها، الگوریتم سیمپلکس نیز وجود دارد که به‌طور گسترده استفاده می‌شود.
اگرچه این الگوریتم در بدترین حالت، زمان چندجمله‌ای ندارد، اما در عمل عملکرد خوبی دارد.
\itm
اینجا وارد جزئیات هیچ کدام از الگوریتم‌ها نمی‌شویم بلکه تنها به بحث چند ایده‌ی مهم می‌پردازیم.
\end{itemframe-s}


\begin{itemframe-s}{الگوریتم‌های برنامه ریزی خطی}{مثالی از یک برنامه خطی دو‌متغیره}
\itm
نخست، مثالی از به‌کارگیری یک روش هندسی برای حل یک برنامه خطی دو‌متغیره ارائه می‌کنیم. هرچند این مثال بلافاصله به یک الگوریتم کارآمد تعمیم نمی‌یابد، اما شما را با مفاهیم مهمی در این زمینه آشنا می‌کند.
\itm
ابتدا برنامه خطی متغیره زیر را در نظر بگیرید:
\begin{align*}
\text{ maximize} & x_1 + x_2\\
\text{ to subject} & 4x_1 - x_2 \geq 8 \\
& 2x_1 + x_2 \geq 10  \\
& 5x_1 - 2x_2 \leq -2 \\
& x_1, x_2 \geq 0
\end{align*}
\end{itemframe-s}


\begin{itemframe-s}{الگوریتم‌های برنامه ریزی خطی}{مثالی از یک برنامه خطی دو‌متغیره}
\itm
اشتراک قیود یک ناحیه محدب (ناحیه آبی رنگ) را تشکیل می‌دهند که همان ناحیه مجاز است:
\pic[.6]{2}
\end{itemframe-s}


\begin{itemframe-s}{الگوریتم‌های برنامه ریزی خطی}{مثالی از یک برنامه خطی دو‌متغیره}
\itm
در دو بعد، می‌توان به کمک یک روش بصری بهینه‌سازی را انجام داد؛
به ازای یک عدد حقیقی مانند
$z$،
مجموعه نقاطی که در آن‌ها
$x_1 + x_2 = z$
برقرار است ، یک خط با شیب $-1$ تشکیل می‌دهند.
برای مثال،
$x_1 + x_2 = 0$
خطی با شیب
$-1$
از مبدأ است.
\itm
اشتراک این خط با ناحیه‌ی مجاز، مجموعه‌ای از نقاط مجاز با مقدار تابع هدف صفر است که در این حالت تنها شامل نقطه $(0,0)$ می‌شود.
\itm
به طور کلی، برای هر مقدار $z$، اشتراک خط
$x_1 + x_2 = z$
با ناحیه‌ی مجاز، مجموعه‌ای از نقاط مجاز با مقدار تابع هدف $z$ خواهد بود.
\itm
از آن‌جا که ناحیه‌ی مجاز در این شکل کراندار است، باید یک مقدار بیشینه $z$ وجود داشته باشد به طوری که اشتراک خط $x_1 + x_2 = z$ با ناحیه‌ی مجاز تهی نباشد.
همانطور که در شکل نشان داده شده، رأس ناحیه‌ی مجاز در
$(2, 6)$
مقدار بهینه را با ارزش تابع هدف $8$ به دست می‌دهد.
\end{itemframe-s}


\begin{itemframe-s}{الگوریتم‌های برنامه ریزی خطی}{مثالی از یک برنامه خطی دو‌متغیره}
\itm
اتفاقی نیست که راه‌حل بهینه در یک رأس ناحیه‌ی مجاز ظاهر شد.
بیشترین مقدار $Z$ برای خط $x_1 + x_2 = z$ الزاماً بر مرز ناحیه‌ی مجاز رخ می‌دهد.
\itm
بنابراین اشتراک این خط با مرز ناحیه‌ی مجاز، یک رأس منفرد یا یک پاره‌خط است.
اگر اشتراک یک رأس باشد، همان رأس تنها راه‌حل بهینه است.
\itm
اگر اشتراک یک پاره‌خط باشد، تمام نقاط آن پاره‌خط مقدار تابع هدف یکسانی دارند، و دو انتهای آن نیز راه‌حل بهینه خواهند بود.
چون هر انتهای پاره‌خط یک رأس است، در این حالت هم یک راه‌حل بهینه در رأس وجود دارد.
\end{itemframe-s}


%todo half space and hyperplane should be known in LA. add it to appendix?
\begin{itemframe-s}{الگوریتم‌های برنامه ریزی خطی}{مثالی از یک برنامه خطی دو‌متغیره}
\itm
اگرچه ترسیم گرافیکی برای مسائل با بیش از دو متغیر آسان نیست، اما همان شهود برقرار می‌ماند. برای سه متغیر، هر قید یک نیم‌فضا در فضای سه‌بعدی تعریف می‌کند.
\itm
اشتراک این نیم‌فضاها ناحیه‌ی مجاز را تشکیل می‌دهد. مجموعه نقاطی که مقدار تابع هدف برابری دارند، اکنون یک صفحه است.
\itm
اگر همه ضرایب تابع هدف نامنفی باشند و مبدأ یک راه‌حل مجاز باشد، با دور کردن این صفحه از مبدأ در جهت عمود بر تابع هدف، نقاطی با مقادیر بزرگ‌تر تابع هدف به دست می‌آید. (در غیر این صورت، تصویر شهودی کمی پیچیده‌تر می‌شود.)
\itm
همانند دو بعد، چون ناحیه‌ی مجاز محدب است، مجموعه نقاطی که مقدار بهینه را می‌سازند، باید شامل یک رأس باشند.
\end{itemframe-s}


\begin{itemframe-s}{الگوریتم‌های برنامه ریزی خطی}{مثالی از یک برنامه خطی دو‌متغیره}
\itm
به طور کلی، اگر $n$ متغیر وجود داشته باشد، هر قید یک نیم‌فضا در فضای $n$ بعدی تعریف می‌کند.
\itm
ناحیه‌ی مجاز حاصل از اشتراک این نیم‌فضاها یک سیمپلکس
\fn{simplex}
 نامیده می‌شود.
\itm
تابع هدف اکنون یک ابرصفحه است و به دلیل محدب بودن، همچنان راه‌حل بهینه در یکی از رأس‌های سیمپلکس رخ می‌دهد.
\end{itemframe-s}


\begin{itemframe-s}{الگوریتم‌های برنامه ریزی خطی}{مثالی از یک برنامه خطی دو‌متغیره}
\itm
الگوریتم سیمپلکس از یک رأس سیمپلکس آغاز کرده و در هر تکرار از یک رأس به رأس مجاور حرکت می‌کند که مقدار تابع هدف در آن بزرگ‌تر یا مساوی با مقدار کنونی باشد.
\itm
الگوریتم زمانی متوقف می‌شود که به یک بیشینه محلی برسد، یعنی رأسی که همه رأس‌های مجاور آن مقدار تابع هدف کوچک‌تری دارند.
ثابت می‌شود، این بهینه‌ی موضعی همان بهینه‌ی سراسری است.
\itm
الگوریتم سیمپلکس اغلب مسائل برنامه‌ریزی خطی را سریع حل می‌کند.
اما برای برخی ورودی‌های خاص که با دقت طراحی شده‌ باشند، زمان اجرای نمایی دارد.
\end{itemframe-s}